\clearpage
\begin{landscape}
\begingroup
\fontsize{10}{13}\selectfont
\fancyhead[L]{\small NACHLASS. STÖRUNGEN DER PALLAS}
\renewcommand{\headrulewidth}{0pt}
\setlength{\lineskip}{4pt}\setlength{\lineskiplimit}{4pt}
\setlength{\arraycolsep}{4pt}
\makeatletter\tagsleft@true\makeatother
\newcommand{\VFivePair}{\mathord{\vcenter{\offinterlineskip\hbox{\rule{.12em}{.12em}}\kern.30em\hbox{\rule{.12em}{.12em}}}}}
\newcommand{\VFivePoint}{\mathord{.}\mkern5mu}
\newcommand{\VFiveTenPoints}{\VFivePoint\VFivePoint\VFivePoint\VFivePoint\VFivePoint\VFivePoint\VFivePoint\VFivePoint\VFivePoint\VFivePoint}
\pdfbookmark[2]{Allgemeine Störungen, Erste Rechnung, [5]}{pallas-V-fragment5}
\SourceLeaf{511}
\begin{center}[5.]\end{center}
\par\noindent\hspace*{\parindent}[Handelt es sich darum, eine Function zweier Argumente von der Form
\[
\begin{aligned}
T&=\gamma_{0,0}+\gamma_{1,0}\cos x+\gamma_{2,0}\cos2x+\cdots\\
 &\quad+\delta_{1,0}\sin x+\delta_{2,0}\sin2x+\cdots\\[4pt]
 &\quad+\gamma_{0,1}\cos y+\gamma_{1,1}\cos(x+y)+\gamma_{2,1}\cos(2x+y)+\cdots\\
 &\quad+\delta_{0,1}\sin y+\delta_{1,1}\sin(x+y)+\delta_{2,1}\sin(2x+y)+\cdots\\[4pt]
 &\qquad\qquad\VFivePair+\epsilon_{1,1}\cos(x-y)+\epsilon_{2,1}\cos(2x-y)+\cdots\\
 &\qquad\qquad+\zeta_{1,1}\sin(x-y)+\zeta_{2,1}\sin(2x-y)+\cdots\\[4pt]
 &\quad+\gamma_{0,2}\cos2y+\gamma_{1,2}\cos(x+2y)+\gamma_{2,2}\cos(2x+2y)+\cdots\\
 &\quad+\delta_{0,2}\sin2y+\delta_{1,2}\sin(x+2y)+\delta_{2,2}\sin(2x+2y)+\cdots\\[4pt]
 &\qquad\qquad+\epsilon_{1,2}\cos(x-2y)+\epsilon_{2,2}\cos(2x-2y)+\cdots\\
 &\qquad\qquad+\zeta_{1,2}\sin(x-2y)+\zeta_{2,2}\sin(2x-2y)+\cdots\\[4pt]
 &\quad+\VFiveTenPoints\VFiveTenPoints\VFiveTenPoints\VFiveTenPoints\VFivePoint\VFivePoint\VFivePoint\VFivePoint\VFivePoint\VFivePoint\VFivePoint
\end{aligned}\tag*{1)}
\]
zu entwickeln, so kann man ebenfalls das in den vorigen Artikeln auseinandergesetzte Verfahren anwenden. Man kann die Function \(T\) in folgender Form schreiben:
\[
\begin{aligned}
T&=\alpha_0+\alpha_1\cos x+\alpha_2\cos2x+\cdots\\
 &\quad+\beta_1\sin x+\beta_2\sin2x+\cdots
\end{aligned}\tag*{2)}
\]
wo die \(\alpha_0,\ \alpha_1,\ \alpha_2\ldots\beta_1,\ \beta_2\ldots\) Functionen von \(y\) allein sind, welche sich wieder in der Form
\[
\begin{aligned}
\alpha_0&=p_{0,0}+p_{0,1}\cos y+p_{0,2}\cos2y+\cdots\\
 &\quad+q_{0,1}\sin y+q_{0,2}\sin2y+\cdots\\[4pt]
\alpha_1&=p_{1,0}+p_{1,1}\cos y+p_{1,2}\cos2y+\cdots\\
 &\quad+q_{1,1}\sin y+\VFivePair q_{1,2}\sin2y+\cdots\\
 &\text{etc.}\\[4pt]
\beta_1&=r_{1,0}+r_{1,1}\cos y+r_{1,2}\cos2y+\cdots\\
 &\quad+s_{1,1}\sin y+s_{1,2}\sin2y+\cdots\\
 &\text{etc.}
\end{aligned}\tag*{3)}
\]
darstellen lassen; die \(p_{0,0},\ p_{0,1},\ldots q_{0,1},\ q_{0,2},\ldots r_{1,0},\ r_{1,1},\ldots s_{1,1},\ s_{1,2},\ldots\) sind constante Coefficienten.

\par\noindent\hspace*{\parindent}Es seien nun für die Werthe \(x=0,\ a,\ 2a,\ldots(\mu-1)a\) und \(y=0,\ a',\ 2a',\ldots(\nu-1)a'\) die zugehörigen Werthe der Function \(T\) gegeben, und zwar sei \(T=A_{0,0}\) für \(x=0,\ y=0\); \(T=A_{1,0}\) für \(x=a,\ y=0\); \(T=A_{2,0}\) für \(x=2a,\ y=0\); \(\ldots T=A_{0,1}\) für \(x=0,\ y=a'\); \(\ldots T=A_{\mu-1,\nu-1}\) für \(x=(\mu-1)a,\ y=(\nu-1)a'\); es seien also im Ganzen \(\mu\nu\) Werthe von \(T\) gegeben; den vorigen Artikeln entsprechend soll auch hier vorausgesetzt werden, dass \(\mu a=\nu a'=360^0\) ist, und dass sowohl \(\mu\) wie \(\nu\) eine gerade Zahl ist. Man bezeichne mit \(T_0\) den Werth, den \(T\) annimmt, wenn \(y=0\) ist, mit \(T_1\) den Werth von \(T\) für \(y=a'\), mit \(T_2\) den für \(y=2a'\), \(\ldots\) mit \(T_{\nu-1}\) den für \(y=(\nu-1)a'\), so dass also \(T_0,\ T_1\) etc. Functionen von \(x\) allein sind, die sich auch wieder in der Form
\SourceLeaf{512}
\[
\begin{aligned}
T_0&=\alpha_{0,0}+\alpha_{1,0}\cos x+\alpha_{2,0}\cos2x+\cdots\\
 &\quad+\beta_{1,0}\sin x+\beta_{2,0}\sin2x+\cdots\\[4pt]
T_1&=\alpha_{0,1}+\alpha_{1,1}\cos x+\alpha_{2,1}\cos2x+\cdots\\
 &\quad+\beta_{1,1}\sin x+\beta_{2,1}\sin2x+\cdots\\[4pt]
 &\VFiveTenPoints\VFiveTenPoints\VFiveTenPoints\VFivePoint\VFivePoint\\[4pt]
T_{\nu-1}&=\alpha_{0,\nu-1}+\alpha_{1,\nu-1}\cos x+\alpha_{2,\nu-1}\cos2x+\cdots\\
 &\quad+\beta_{1,\nu-1}\sin x+\beta_{2,\nu-1}\sin2x+\cdots
\end{aligned}\tag*{4)}
\]
darstellen.

\par\noindent\hspace*{\parindent}Offenbar sind nun \(A_{0,0},\ A_{1,0},\ A_{2,0}\ldots A_{\mu-1,0}\) der Reihe nach die Werthe, welche \(T_0\) annimmt für \(x=0,\ a,\ 2a,\ldots(\mu-1)a\); es lassen sich aus ihnen genau nach den Vorschriften der Artikel [2]---[4] die Coefficienten \(\alpha_{0,0},\ \alpha_{1,0},\ldots\beta_{1,0},\ \beta_{2,0}\ldots\) bestimmen; z.~B. hat man:
\[
\begin{aligned}
\mu\alpha_{0,0}&=A_{0,0}+A_{1,0}+A_{2,0}+\cdots A_{\mu-1,0}\\[4pt]
\frac\mu2\alpha_{1,0}&=A_{0,0}+A_{1,0}\cos a+A_{2,0}\cos2a+\cdots+A_{\mu-1,0}\cos(\mu-1)a\\[4pt]
\frac\mu2\beta_{1,0}&=A_{0,0}+A_{1,0}\sin a+A_{2,0}\sin2a+\cdots+A_{\mu-1,0}\sin(\mu-1)a\\
 &\text{etc.}
\end{aligned}
\]

\par\noindent\hspace*{\parindent}Ebenso sind \(A_{0,1},\ A_{1,1},\ A_{2,1}\ldots A_{\mu-1,1}\) der Reihe nach die Werthe von \(T_1\) für \(x=0,\ a,\ 2a,\ldots(\mu-1)a\), und aus ihnen lassen sich die Coefficienten \(\alpha_{0,1},\ \alpha_{1,1},\ \alpha_{2,1}\ldots\beta_{1,1},\ \beta_{2,1}\ldots\) in derselben Weise bestimmen. Dasselbe Verfahren ist anzuwenden für \(T_2,\ T_3,\ldots T_{\nu-1}\).

\par\noindent\hspace*{\parindent}Nach Ausführung dieser Rechnungen hat man die Ausdrücke 4 für die Functionen \(T_0,\ T_1,\ldots T_{\nu-1}\) vollkommen bestimmt.

\par\noindent\hspace*{\parindent}Des Weiteren sind nun aber die jetzt bekannten Coefficienten \(\alpha_{0,0},\ \alpha_{0,1},\ \alpha_{0,2}\ldots\alpha_{0,\nu-1}\) der Reihe nach die Werthe, welche \(\alpha_0\) annimmt für \(y=0,\ a',\ 2a'\ldots(\nu-1)a'\) und \(\alpha_{1,0},\ \alpha_{1,1},\ \alpha_{1,2},\ldots\alpha_{1,\nu-1}\) die Werthe von \(\alpha_1\) für \(y=0,\ a',\ 2a'\ldots(\nu-1)a'\) etc.; ebenso \(\beta_{1,0},\ \beta_{1,1},\ \beta_{1,2},\ldots\beta_{1,\nu-1}\) die Werthe von \(\beta_1\) für die genannten Werthe von \(y\). Es lassen sich also wieder nach den vorigen Artikeln
\[
\begin{array}{@{}ll@{\quad}l@{\quad}l@{\quad}l@{}}
\text{die}&\text{Coefficienten}&p_{0,0},\ p_{0,1},\ p_{0,2},\ldots q_{0,1},\ q_{0,2}\ldots&\text{aus}&\alpha_{0,0},\ \alpha_{0,1},\ \alpha_{0,2}\ldots\\[4pt]
\text{»}&\text{»}&p_{1,0},\ p_{1,1},\ p_{1,2},\ldots q_{1,1},\ q_{1,2},\ldots&\text{aus}&\alpha_{1,0},\ \alpha_{1,1},\ \alpha_{1,2}\ldots\\[4pt]
\multicolumn{5}{@{}c@{}}{\VFiveTenPoints\VFiveTenPoints\VFiveTenPoints\VFiveTenPoints\VFiveTenPoints\VFivePoint\VFivePoint\VFivePoint\VFivePoint\VFivePoint}\\[4pt]
\text{»}&\text{»}&r_{1\,0},\ r_{1,1},\ r_{1,2},\ldots s_{1,1},\ s_{1,2},\ldots&\text{aus}&\beta_{1,0},\ \beta_{1,1},\ \beta_{1,2}\ldots\\
\multicolumn{5}{c}{\text{etc.}}
\end{array}
\]
bestimmen.

\par\noindent\hspace*{\parindent}Hiemit sind auch die Entwickelungen 3 für die \(\alpha_0,\ \alpha_1,\ \alpha_2,\ldots\beta_1,\ \beta_2,\ldots\) vollkommen bekannt, und indem man diese Werthe in die Reihe 2 für \(T\) einsetzt, erhält man die vollständige Entwickelung der Function \(T\) nach beiden Veränderlichen, wobei man die auftretenden Producte trigonometrischer Functionen in einfache Sinus und Cosinus der Summen und Differenzen umwandelt, um \(T\) in der Form 1 zu erhalten.

\par\noindent\hspace*{\parindent}Es mögen noch die Formeln zur Berechnung von \(\gamma_{0,0},\allowbreak\ \gamma_{1,0},\allowbreak\ \gamma_{2,0},\allowbreak\ldots\delta_{1,0},\allowbreak\ \delta_{2,0},\allowbreak\ldots\epsilon_{1,1},\allowbreak\ \epsilon_{2,1},\allowbreak\ldots\zeta_{1,1},\allowbreak\ \zeta_{2,1},\allowbreak\ldots\) aus \(p_{0,0},\allowbreak\ p_{0,1},\allowbreak\ldots q_{0,1},\allowbreak\ q_{0,2},\allowbreak\ldots r_{1,0},\allowbreak\ r_{1,1},\allowbreak\ldots s_{1,1},\allowbreak\ s_{1,2},\allowbreak\ldots\) angeführt werden, obwohl man die Rechnung eben so bequem rein mechanisch ausführt:
\[
\begin{aligned}
\gamma_{0,0}&=p_{0,0},&\gamma_{0,1}&=p_{0,1},&\gamma_{0,2}&=p_{0,2}\quad\text{etc.}\\
&&\delta_{0,1}&=q_{0,1},&\delta_{0,2}&=q_{0,2}\quad\text{etc.}
\end{aligned}
\]
\SourceLeaf{513}
\[
\begin{aligned}
\gamma_{1,0}&=p_{1,0},&\gamma_{1,1}&=\tfrac12(p_{1,1}-s_{1,1}),&\gamma_{1,2}&=\tfrac12(p_{1,2}-s_{1,2})\quad\text{etc.}\\
\delta_{1,0}&=r_{1,0},&\delta_{1,1}&=\tfrac12(r_{1,1}+q_{1,1}),&\delta_{1,2}&=\tfrac12(r_{1,2}+q_{1,2})\quad\text{etc.}\\[8pt]
&&\epsilon_{1,1}&=\tfrac12(p_{1,1}+s_{1,1}),&\epsilon_{1,2}&=\tfrac12(p_{1,2}+s_{1,2})\quad\text{etc.}\\
&&\zeta_{1,1}&=\tfrac12(r_{1,1}-q_{1,1}),&\zeta_{1,2}&=\tfrac12(r_{1,2}-q_{1,2})\quad\text{etc.}\\[8pt]
\gamma_{2,0}&=p_{2,0},&\gamma_{2,1}&=\tfrac12(p_{2,1}-s_{2,1}),&\gamma_{2,2}&=\tfrac12(p_{2,2}-s_{2,2})\quad\text{etc.}\\
\delta_{2,0}&=r_{2,0},&\delta_{2,1}&=\tfrac12(r_{2,1}+q_{2\,1}),&\delta_{2,2}&=\tfrac12(r_{2,2}+q_{2,2})\quad\text{etc.}\\[8pt]
&&\epsilon_{2,1}&=\tfrac12(p_{2,1}+s_{2,1}),&\epsilon_{2,2}&=\tfrac12(p_{2,2}+s_{2,2})\quad\text{etc.}\\
&&\zeta_{2,1}&=\tfrac12(r_{2,1}-q_{2,1}),&\zeta_{2,2}&=\tfrac12(r_{2,2}-q_{2,2})\quad\text{etc.}
\end{aligned}
\]
\begin{center}u.~s.~w.\end{center}
\par\noindent\hspace*{\parindent}Man vergleiche zu diesem Artikel die Artikel 25---27 der Theoria interpolationis (Band III Seite 303).]
\clearpage
\endgroup
\end{landscape}
